\section{Ontologies and the Semantic Web}
\label{sec:ontologies}
\slidesrc{Ontologies: 1--48}

Lexical semantics (Section~\ref{sec:lexical-semantics}) organises the meanings of
\emph{words}; WordNet already looks like a hierarchy of concepts. \textbf{Ontologies} go
one step further: they describe the \emph{concepts of a domain} (wines, courses,
companies, \ldots), their properties and relations, in a formal language that a machine
can reason with. The \textbf{Semantic Web} is the project of using ontologies to give the
content of the Web a well-defined, machine-readable meaning.

This section is a soft introduction to both. It covers:
\begin{enumerate}[nosep]
  \item what an ontology is, why semantic interoperability matters, how ontologies
        differ from databases, and knowledge graphs (Section~\ref{sec:what-is-an-ontology});
  \item a methodology for building an ontology, and the Protégé editor
        (Section~\ref{sec:building-ontologies});
  \item the ontology languages RDF, RDF Schema and OWL
        (Section~\ref{sec:ontology-languages});
  \item the Semantic Web, its layered ``stack'' and what came after it
        (Section~\ref{sec:semantic-web}).
\end{enumerate}

\subsection{What Is an Ontology}
\label{sec:what-is-an-ontology}
\slidesrc{Ontologies: 3--7, 10--21}

\subsubsection{Two definitions}
\slidesrc{Ontologies: 3--7}

The word comes from philosophy, where \emph{ontology} is the study of being. Computer
science borrowed it in the early 1990s with a much more practical meaning:
\begin{itemize}[nosep]
  \item ``The study of that which is common to all things which exist''
        (Aristotle, \emph{Metaphysics});
  \item ``\textbf{Explicit specification of a conceptualization}'' (Tom Gruber,
        \emph{A Translation Approach to Portable Ontology Specifications}, Knowledge
        Acquisition 5(2):199--220, 1993).
\end{itemize}

\begin{nbintuition}[reading Gruber's definition]
A \emph{conceptualization} is the way we carve up a domain: which kinds of things exist
(wines, producers, foods), how they are related (a wine \emph{is produced by} a
producer), which constraints hold. It normally lives only in people's heads.
\emph{Explicit specification} means writing it down, in a formal language, so that it
can be shared between people and programs and checked by a machine.
\end{nbintuition}

\subsubsection{Motivation: semantic interoperability}
\slidesrc{Ontologies: 10--12}

\begin{nbdef}{Semantic interoperability}
\textbf{Semantic interoperability} is ``the ability of computer systems to exchange data
with unambiguous, shared meaning''.
\end{nbdef}

Two systems can exchange bytes perfectly and still disagree on what they mean: is
\texttt{price} with or without VAT, is \texttt{date} the order date or the shipping
date? Semantic interoperability has been one of the most active research fields of the
last twenty-five years: the advent of ontologies in the early nineties, Web Services at
the beginning of the new millennium, and their combination into Semantic Web Services
all aim at making the ``semantic interoperability dream'' come true.

\begin{nbwhy}
The interest is economic as much as scientific. Walker et al.\ (\emph{The Value of
Healthcare Information Exchange and Interoperability}, Health Affairs, 2005) estimated
that a fully standardised system for information exchange and interoperability in US
health care could yield a net value of \textbf{77.8 billion dollars per year} once
implemented. More recent work (e.g.\ on semantic interoperability in health care, 2018,
and the European Data Conference on Reference Data and Semantics, 2023) confirms the
expected advantages.
\end{nbwhy}

\subsubsection{Gruber's definition, 2008}
\slidesrc{Ontologies: 13--15}

Fifteen years after his first definition, Gruber gave a more detailed one (\emph{Ontology},
Encyclopedia of Database Systems, Springer, 2008).

\begin{nbdef}{Ontology (Gruber, 2008)}
In computer and information science, an \textbf{ontology} defines a set of
\textbf{representational primitives} with which to model a domain of knowledge or
discourse. The primitives are typically
\begin{itemize}[nosep]
  \item \textbf{classes} (or sets),
  \item \textbf{attributes} (or properties),
  \item \textbf{relationships} (relations among class members).
\end{itemize}
Their definitions include information about their \emph{meaning} and \emph{constraints}
on their logically consistent application.
\end{nbdef}

From the point of view of databases:
\begin{itemize}[nosep]
  \item an ontology is a level of abstraction of data models, analogous to hierarchical
        and relational models, but intended for modelling \emph{knowledge} about
        individuals, their attributes and their relationships to other individuals;
  \item ontology languages abstract away from data structures and implementation
        strategies, and are closer in expressive power to \textbf{first-order logic} than
        database languages. Ontologies are therefore said to be at the
        \textbf{``semantic'' level}, database schemas at the ``logical'' or ``physical''
        level;
  \item being independent of the lower-level data models, ontologies are used to
        integrate heterogeneous databases, to make disparate systems interoperate and to
        specify interfaces to knowledge-based services. In the Semantic Web standards
        they are an explicit layer (Section~\ref{sec:semantic-web}), with standard
        languages and many commercial and open source tools.
\end{itemize}

\subsubsection{An example: a wine ontology}
\slidesrc{Ontologies: 16}

The slides show a small ontology of wines and foods (in Italian). Each node is a class;
each child is a subclass (``is a kind of'') of its parent, exactly like the hypernym
hierarchy of WordNet (Section~\ref{sec:wordnet}).

\begin{center}
\begin{tikzpicture}[
    font=\small, x=0.6cm, y=-0.48cm,
    cls/.style={anchor=west, inner sep=2pt},
    inner/.style={cls, text=nbblue!80!black},
    edge/.style={draw=nbgray, thick}]
  \node[inner] (thing) at (0,0) {Thing};
  \node[inner] (vino) at (1,1) {Vino \emph{(wine)}};
  \node[inner] (vb) at (2,2) {Vino Bianco \emph{(white wine)}};
  \node[cls] (pb) at (3,3) {Pinot Bianco};
  \node[inner] (ri) at (3,4) {Riesling};
  \node[cls] (rs) at (4,5) {Riesling Secco \emph{(dry)}};
  \node[cls] (rd) at (4,6) {Riesling Dolce \emph{(sweet)}};
  \node[cls] (ch) at (3,7) {Chardonnay};
  \node[cls] (vr) at (2,8) {Vino Rosso \emph{(red)}};
  \node[cls] (vro) at (2,9) {Vino Rosè \emph{(rosé)}};
  \node[cls] (vd) at (2,10) {Vino da dessert \emph{(dessert wine)}};
  \node[inner] (cibo) at (1,11) {Cibo \emph{(food)}};
  \node[inner] (carne) at (2,12) {Carne \emph{(meat)}};
  \node[cls] (cb) at (3,13) {Carni Bianche \emph{(white meat)}};
  \node[cls] (cr) at (3,14) {Carni Rosse \emph{(red meat)}};
  \node[inner] (pesce) at (2,15) {Pesce \emph{(fish)}};
  \node[cls] (ps) at (3,16) {Pesce d'acqua salata \emph{(saltwater)}};
  \node[cls] (pd) at (3,17) {Pesce d'acqua dolce \emph{(freshwater)}};
  \node[cls] (prod) at (1,18) {Produttore \emph{(producer)}};
  \draw[edge] ([xshift=0.3cm]thing.south west) |- (vino.west);
  \draw[edge] ([xshift=0.3cm]vino.south west) |- (vb.west);
  \draw[edge] ([xshift=0.3cm]vb.south west) |- (pb.west);
  \draw[edge] ([xshift=0.3cm]vb.south west) |- (ri.west);
  \draw[edge] ([xshift=0.3cm]ri.south west) |- (rs.west);
  \draw[edge] ([xshift=0.3cm]ri.south west) |- (rd.west);
  \draw[edge] ([xshift=0.3cm]vb.south west) |- (ch.west);
  \draw[edge] ([xshift=0.3cm]vino.south west) |- (vr.west);
  \draw[edge] ([xshift=0.3cm]vino.south west) |- (vro.west);
  \draw[edge] ([xshift=0.3cm]vino.south west) |- (vd.west);
  \draw[edge] ([xshift=0.3cm]thing.south west) |- (cibo.west);
  \draw[edge] ([xshift=0.3cm]cibo.south west) |- (carne.west);
  \draw[edge] ([xshift=0.3cm]carne.south west) |- (cb.west);
  \draw[edge] ([xshift=0.3cm]carne.south west) |- (cr.west);
  \draw[edge] ([xshift=0.3cm]cibo.south west) |- (pesce.west);
  \draw[edge] ([xshift=0.3cm]pesce.south west) |- (ps.west);
  \draw[edge] ([xshift=0.3cm]pesce.south west) |- (pd.west);
  \draw[edge] ([xshift=0.3cm]thing.south west) |- (prod.west);
\end{tikzpicture}
\captionof{figure}{The wine ontology of the slides as a class hierarchy (translations in
italics).}
\label{fig:wine-ontology}
\end{center}

The slide also labels the \textbf{levels} of the hierarchy: \emph{top level} (the most
general classes, such as \emph{Vino}), \emph{middle level} (\emph{Vino Bianco},
\emph{Riesling}, \emph{Chardonnay}, \ldots) and \emph{bottom level} (the most specific
classes, \emph{Riesling Secco} and \emph{Riesling Dolce}). A full ontology would also
have \emph{properties} linking the hierarchies, e.g.\ a wine \emph{is made by} a
\emph{Produttore}, a wine \emph{goes well with} a \emph{Cibo}.

\subsubsection{Ontologies and databases}
\slidesrc{Ontologies: 17--18}

\begin{center}
\small
\begin{tabularx}{\linewidth}{@{}l >{\raggedright\arraybackslash}X >{\raggedright\arraybackslash}X@{}}
\toprule
 & \textbf{Ontologies} & \textbf{Databases} \\
\midrule
Purpose & sharing, reuse, defining semantics & storing large amounts of structured
  data \\
Expressive power & close to first-order logic; description logics & not as powerful as
  first-order logic \\
Query standards & no consolidated standard, few proposals (e.g.\ SPARQL) & consolidated
  standards (e.g.\ SQL) \\
Management systems & Protégé, Jena (not commercial) & many solid commercial tools and
  environments (Oracle, \ldots) \\
\bottomrule
\end{tabularx}
\captionof{table}{Ontologies and databases compared.}
\end{center}

\begin{nbwarning}[SPARQL today]
SPARQL has since become a W3C Recommendation (2008; SPARQL 1.1 in 2013), so it is now the
standard query language for RDF data, even if it is far less widespread than SQL. The
comparison in the slides reflects an earlier state of the field.
\end{nbwarning}

\subsubsection{Knowledge graphs}
\slidesrc{Ontologies: 19--21}

\begin{nbdef}{Knowledge graph}
``The heart of the \textbf{knowledge graph} is a knowledge model: a collection of
interlinked descriptions of concepts, entities, relationships and events. Knowledge
graphs put data in context via linking and semantic metadata and this way provide a
framework for data integration, unification, analytics and sharing'' (Ontotext).
\end{nbdef}

\begin{center}
\begin{tikzpicture}[
    font=\scriptsize,
    ent/.style={circle, draw=none, fill=nbteal!80, text=white, minimum size=1.25cm,
                align=center, font=\scriptsize\bfseries},
    co/.style={ent, fill=nbgray!80},
    typ/.style={ent, fill=nborange!90},
    rel/.style={-{Stealth}, thick, nbgray},
    inf/.style={-{Stealth}, thick, dashed, nbred!80},
    lab/.style={midway, fill=white, inner sep=1pt, font=\tiny}]
  \node[co]  (big)  at (0,  4.4) {Big\\Bucks\\Café};
  \node[co]  (glob) at (0,  2.2) {Global\\Invest.\\Inc.};
  \node[co]  (my)   at (0,  0)   {My\\Local\\Café};
  \node[ent] (sea)  at (3.4,4.4) {Seattle};
  \node[ent] (wb)   at (3.4,2.2) {West\\Bay};
  \node[ent] (nyc)  at (3.4,0)   {New York\\City};
  \node[ent] (wa)   at (6.8,4.4) {Washing-\\ton};
  \node[ent] (cay)  at (6.8,2.2) {Cayman\\Island};
  \node[ent] (ny)   at (6.8,0)   {New\\York};
  \node[ent] (usa)  at (10.2,2.2) {USA};
  \node[typ] (ctry) at (10.2,4.4) {Country};
  \node[typ] (off)  at (10.2,0)   {Offshore\\Zone};
  \draw[rel] (big)  -- node[lab]{controls}     (glob);
  \draw[rel] (glob) -- node[lab]{controls}     (my);
  \draw[rel] (big)  -- node[lab]{located in}   (sea);
  \draw[rel] (glob) -- node[lab]{located in}   (wb);
  \draw[rel] (my)   -- node[lab]{located in}   (nyc);
  \draw[rel] (sea)  -- node[lab]{subregion of} (wa);
  \draw[rel] (wb)   -- node[lab]{subregion of} (cay);
  \draw[rel] (nyc)  -- node[lab]{subregion of} (ny);
  \draw[rel] (wa)   -- node[lab, pos=0.35]{subregion of} (usa);
  \draw[rel] (ny)   -- node[lab, pos=0.22]{subregion of} (usa);
  \draw[rel] (usa)  -- node[lab]{type}         (ctry);
  \draw[rel] (cay)  -- node[lab, pos=0.8]{type} (off);
  \draw[inf] (big.north) to[bend left=25] node[lab]{located in} (wa.north);
  \draw[inf] (my.south)  to[bend right=25] node[lab]{located in} (ny.south);
  \draw[inf] (glob) to[bend left=22] node[lab]{located in} (cay);
\end{tikzpicture}
\captionof{figure}{A knowledge graph of companies and places, after the slides. Solid
arrows are stated facts, dashed arrows are \emph{inferred} facts.}
\label{fig:knowledge-graph}
\end{center}

\begin{nbexample}{what the graph lets us infer}
The stated facts say that \emph{My Local Café} is located in \emph{New York City}, which
is a subregion of \emph{New York}: since \emph{located in} propagates along
\emph{subregion of}, the café is also located in New York, and in the USA. Following the
\emph{controls} links upwards, \emph{My Local Café} is controlled by \emph{Global
Investment Inc.}, located in West Bay, a subregion of the Cayman Islands, which are an
\emph{Offshore Zone}. The slide's conclusion, ``discover suspicious relationships'', is
a query over the graph: \emph{a company controlled by another company from the same
country through an offshore zone}. A news item saying that ``My Local Café is a network
of independent coffee shops'' can be checked against it.
\end{nbexample}

\begin{nbintuition}[ontology vs.\ knowledge graph]
The two terms overlap heavily. A useful way to keep them apart: the \emph{ontology} is
the schema (classes such as \emph{Company}, \emph{Region}; properties such as
\emph{controls}, \emph{located in}; rules such as ``\emph{located in} is transitive
through \emph{subregion of}''), while the \emph{knowledge graph} is the large network of
concrete entities and facts described with that schema.
\end{nbintuition}

\subsection{Building Ontologies: Methodology and Tools}
\label{sec:building-ontologies}
\slidesrc{Ontologies: 22--23, 38}

\subsubsection{Ontology Development 101}
\slidesrc{Ontologies: 22--23}

A widely used methodology is the one of Natalya F.\ Noy and Deborah L.\ McGuinness,
\emph{Ontology Development 101: A Guide to Creating Your First Ontology} (Stanford
Knowledge Systems Laboratory, 2001). It has seven steps:
\begin{enumerate}[nosep]
  \item determine the \textbf{domain and scope} of the ontology;
  \item consider \textbf{reusing} existing ontologies;
  \item enumerate the important \textbf{terms} in the ontology;
  \item define the \textbf{classes} and the \textbf{class hierarchy};
  \item define the \textbf{properties} of classes (\emph{slots});
  \item define the \textbf{facets} of the slots;
  \item create \textbf{instances}.
\end{enumerate}

\begin{nbexample}{the seven steps on the wine ontology}
The guide itself uses wines and foods as its running example, the same domain as
Figure~\ref{fig:wine-ontology}.
\begin{enumerate}[nosep]
  \item \emph{Domain and scope}: wines and foods, to recommend a wine for a meal. Useful
        \emph{competency questions}: ``which wine goes well with fish?'', ``is
        Chardonnay a white wine?''.
  \item \emph{Reuse}: is there already a food or geography ontology we can import?
  \item \emph{Terms}: wine, white wine, Riesling, producer, colour, sugar, region,
        fish, red meat, \ldots
  \item \emph{Classes and hierarchy}: \emph{Riesling} is a subclass of \emph{White
        wine}, which is a subclass of \emph{Wine}.
  \item \emph{Slots}: a wine has a \emph{colour}, a \emph{sugar} level, a \emph{maker}
        (a \emph{Producer}), and \emph{goes well with} some \emph{Food}.
  \item \emph{Facets} (constraints on slots): \emph{colour} takes one value among
        \{red, white, rosé\}; \emph{maker} has exactly one value, of class
        \emph{Producer}.
  \item \emph{Instances}: a specific bottle, e.g.\ a given producer's Riesling of a
        given year, as an individual of class \emph{Riesling}.
\end{enumerate}
\end{nbexample}

\begin{nbintuition}[terms first, classes later]
Steps 3--4 mirror what we did with words in lexical semantics: first collect the
vocabulary of the domain, then decide which terms denote classes, which ones denote
properties, and arrange the classes in an IS-A hierarchy (the hyponymy relation of
Section~\ref{sec:lexical-relations}). The process is iterative: the first version of an
ontology is almost never the final one.
\end{nbintuition}

\subsubsection{Tools: Protégé}
\slidesrc{Ontologies: 38}

\textbf{Protégé} (\url{https://protege.stanford.edu/}) is a free, open source ontology
editor and knowledge-base framework developed at Stanford University. It is written in
Java, is extensible, and provides a plug-and-play environment that makes it a flexible
base for rapid prototyping and application development. It supports two main ways of
modelling ontologies:
\begin{itemize}[nosep]
  \item the \textbf{Protégé-Frames} editor, to build and populate \emph{frame-based}
        ontologies (classes with slots and facets, as in the methodology above);
  \item the \textbf{Protégé-OWL} editor, to build ontologies for the Semantic Web, in
        the W3C's OWL language (Section~\ref{sec:ontology-languages}).
\end{itemize}

\subsection{Ontology Languages: RDF, RDF Schema, OWL}
\label{sec:ontology-languages}
\slidesrc{Ontologies: 24--37}

\subsubsection{Evolution of the languages}
\slidesrc{Ontologies: 25}

Ontology languages for the Web were built on top of each other, starting from the markup
languages of the Web itself.

\begin{center}
\begin{tikzpicture}[
    lay/.style={draw=nbblue!70!black, fill=nbblue!8, minimum height=0.6cm,
                font=\small, inner sep=0pt}]
  \node[lay, minimum width=1.4cm, anchor=south west] at (0,0)    {HTML};
  \node[lay, minimum width=6.6cm, anchor=south west] at (1.4,0)  {XML};
  \node[lay, minimum width=1.6cm, anchor=south west] at (0.7,0.6) {XHTML};
  \node[lay, minimum width=5.4cm, anchor=south west] at (2.6,0.6) {RDF};
  \node[lay, minimum width=4.4cm, anchor=south west] at (3.1,1.2) {RDF Schema};
  \node[lay, minimum width=1.4cm, anchor=south west] at (3.1,1.8) {OWL};
  \node[lay, minimum width=1.4cm, anchor=south west] at (4.6,1.8) {DAML};
  \node[lay, minimum width=1.4cm, anchor=south west] at (6.1,1.8) {OIL};
  \node[lay, minimum width=2.9cm, anchor=south west] at (4.6,2.4) {DAML+OIL};
\end{tikzpicture}
\captionof{figure}{Evolution of the Web and ontology languages: each layer builds on the
ones below it.}
\label{fig:languages-evolution}
\end{center}

\begin{itemize}[nosep]
  \item \textbf{XML} gives a common \emph{syntax} for structured documents, but says
        nothing about what the tags mean.
  \item \textbf{RDF} adds a simple data model (triples, below) on top of XML.
  \item \textbf{RDF Schema} adds a vocabulary for classes, subclasses and properties.
  \item \textbf{DAML} (US) and \textbf{OIL} (Europe) were two richer ontology languages
        built on RDF Schema, merged into \textbf{DAML+OIL}, which in turn was the basis of
        the W3C standard \textbf{OWL}.
\end{itemize}

\subsubsection{RDF: resources, properties, statements}
\slidesrc{Ontologies: 26--29}

\textbf{RDF} (Resource Description Framework) is a W3C Recommendation. It provides a
general, flexible method to decompose any knowledge into small pieces, called
\textbf{triples}, with some rules about the meaning of those pieces.

\begin{nbdef}{RDF data model}
The RDF data model has three object types:
\begin{itemize}[nosep]
  \item a \textbf{resource}: anything that can have a URI (Web identifier). It can be a
        whole Web page, a part of it, a collection of pages, or an object that is not
        directly accessible via the Web (a person, a book);
  \item a \textbf{property}: a specific aspect, characteristic, attribute or relation
        used to describe a resource (\emph{author}, \emph{homepage}). Properties are
        resources too;
  \item a \textbf{statement}: a triple made of a resource, a property and a property
        value (which can be a literal or another resource). The three parts are called
        the \textbf{subject}, \textbf{predicate} and \textbf{object} of the statement.
\end{itemize}
\end{nbdef}

A triple $(O, A, V)$ is more often written as $A(O, V)$, and means: \emph{object $O$ has
attribute $A$ with value $V$}. A set of triples is a directed, labelled \textbf{graph}:
subjects and objects are nodes, each predicate is an edge.

\begin{center}
\begin{tikzpicture}[
    font=\small,
    res/.style={ellipse, draw, align=center, font=\scriptsize\ttfamily, inner sep=1pt},
    lit/.style={rectangle, draw, font=\small, minimum height=0.6cm},
    arr/.style={-{Stealth}, thick},
    lab/.style={font=\scriptsize\ttfamily, fill=white, inner sep=1pt}]
  \node[res] (emp)  at (0,0)    {http://www.w3.org/\\employee/id1321};
  \node[res] (book) at (6.2,0)  {http://www.books.org/\\ISBN0012515866};
  \node[lit] (name) at (2.6,2.0) {Jim Lerners};
  \node[lit] (pr)   at (10.6,0) {\$62};
  \draw[arr] (emp)  -- node[lab, above]{s:authorOf}   (book);
  \draw[arr] (emp)  -- node[lab, left=2pt]{s:hasName} (name);
  \draw[arr] (book) -- node[lab, above]{s:hasPrice}   (pr);
\end{tikzpicture}
\captionof{figure}{An RDF graph with three statements. Ellipses are resources
(identified by URIs), boxes are literal values.}
\label{fig:rdf-graph}
\end{center}

The graph contains three triples:
\begin{align*}
  &\mathit{hasName}(\mathit{id1321}, \text{``Jim Lerners''}),\\
  &\mathit{authorOf}(\mathit{id1321}, \mathit{ISBN0012515866}),\\
  &\mathit{hasPrice}(\mathit{ISBN0012515866}, \text{\$62}).
\end{align*}
Because the book is a
resource, it can be the object of one statement and the subject of another: this is
what makes triples link into a graph.

\begin{nbexample}{RDF in XML syntax}
\begin{minted}{xml}
<?xml version="1.0"?>
<RDF>
  <Description about="http://www.w3schools.com/rdf">
    <author>Jan Egil Refsnes</author>
    <homepage>http://www.w3schools.com</homepage>
  </Description>
</RDF>
\end{minted}
The \texttt{Description} element describes one resource, and each child element is a
property with its value. The document contains two statements:
\begin{center}
\small
\begin{tabular}{@{}lll@{}}
\toprule
\textbf{Subject} & \textbf{Predicate} & \textbf{Object} \\
\midrule
\url{http://www.w3schools.com/rdf} & author   & Jan Egil Refsnes \\
\url{http://www.w3schools.com/rdf} & homepage & \url{http://www.w3schools.com} \\
\bottomrule
\end{tabular}
\end{center}
that is, ``the author of \url{http://www.w3schools.com/rdf} is Jan Egil Refsnes'' and
``the homepage of \url{http://www.w3schools.com/rdf} is
\url{http://www.w3schools.com}''. In the second statement the value is itself a
resource.
\end{nbexample}

\begin{nbwarning}[the example is simplified]
In real RDF/XML the RDF vocabulary lives in its own namespace and the properties in
another one: the root is \texttt{rdf:RDF} with an \texttt{xmlns:rdf} declaration, the
resource is \texttt{rdf:Description rdf:about="\ldots"}, and \texttt{author},
\texttt{homepage} carry a prefix of a vocabulary of our choice. The unprefixed version
above is easier to read but is not valid RDF.
\end{nbwarning}

\subsubsection{RDF Schema}
\slidesrc{Ontologies: 30--31}

RDF alone lets us state facts but not define a vocabulary. \textbf{RDF Schema} (RDFS) is
used to define the structure of the metadata that describe Web resources:
\begin{itemize}[nosep]
  \item it consists of some basic classes and properties, and can be extended to fit
        possibly any domain;
  \item \textbf{classes} are arranged hierarchically (\texttt{rdfs:subClassOf}). The
        root of the hierarchy is \texttt{rdfs:Resource}; \texttt{rdfs:Class} is a
        subclass of \texttt{rdfs:Resource};
  \item \textbf{properties} are instances of the class \texttt{rdf:Property}; they
        describe resources by assigning values to them, and are resources themselves.
        Their use can be constrained to members of certain classes:
        \texttt{rdfs:domain} gives the class of the subjects, \texttt{rdfs:range} the
        class of the values.
\end{itemize}

\begin{nbexample}{teachers and courses in RDFS}
\begin{minted}{xml}
<rdfs:Class rdf:ID="Teacher">
  <rdfs:comment>Teacher Class</rdfs:comment>
  <rdfs:subClassOf rdf:resource="#Person"/>
</rdfs:Class>
<rdfs:Class rdf:ID="Course">
  <rdfs:comment>Course Class</rdfs:comment>
  <rdfs:subClassOf
  rdf:resource="http://www.w3.org/1999/02/22-rdf-syntax-ns#Resource"/>
</rdfs:Class>
<rdf:Property rdf:ID="teacher">
  <rdfs:comment>Teacher of a course</rdfs:comment>
  <rdfs:domain rdf:resource="#Course"/>
  <rdfs:range rdf:resource="#Teacher"/>
</rdf:Property>
<rdf:Property rdf:ID="name">
  <rdfs:comment>Name of a Person or Course</rdfs:comment>
  <rdfs:domain rdf:resource="#Person"/>
  <rdfs:domain rdf:resource="#Course"/>
  <rdfs:range
  rdf:resource="http://www.w3.org/1999/02/22-rdf-syntax-ns#Literal"/>
</rdf:Property>
\end{minted}
\emph{Teacher} is a subclass of \emph{Person}; \emph{teacher} is a property linking a
\emph{Course} (domain) to a \emph{Teacher} (range): ``the teacher of this course is
\ldots''. \emph{name} gives a literal (a string) to persons and courses.
\end{nbexample}

\begin{nbwarning}[two domains mean ``and'', not ``or'']
In RDFS semantics, several \texttt{rdfs:domain} statements on the same property are
\emph{all} applied: whatever has a \emph{name} is inferred to be \emph{both} a
\emph{Person} and a \emph{Course}. To express ``a Person \emph{or} a Course'', the
domain must be a single class that is the union of the two (OWL can define it with
\texttt{owl:unionOf}), or a common superclass.
\end{nbwarning}

\subsubsection{OWL: the Web Ontology Language}
\slidesrc{Ontologies: 32--37}

\textbf{OWL} extends RDF Schema: it uses the RDF meaning of classes and properties
(\texttt{rdfs:Class}, \texttt{rdfs:subClassOf}, \ldots) and adds language primitives for
a richer expressiveness. It comes in three species, which trade expressive power for
ease of reasoning:

\begin{center}
\small
\begin{tabularx}{\linewidth}{@{}l >{\raggedright\arraybackslash}X >{\raggedright\arraybackslash}X >{\raggedright\arraybackslash}X@{}}
\toprule
\textbf{Species} & \textbf{What it is} & \textbf{Advantage} & \textbf{Disadvantage} \\
\midrule
OWL Full & all OWL primitives, combinable in arbitrary ways with RDF and RDF Schema &
  fully upward compatible with RDF & so powerful that it is \textbf{undecidable} \\
OWL DL & a sublanguage of OWL Full that restricts how OWL and RDF constructors can be
  used (DL = description logic) & permits \textbf{efficient reasoning} support & loses
  full compatibility with RDF \\
OWL Lite & a further restriction of OWL DL to a subset of the constructors & easier to
  grasp (users) and to implement (tool builders) & restricted expressivity \\
\bottomrule
\end{tabularx}
\captionof{table}{The three species of OWL, from the most to the least expressive.}
\end{center}

\begin{nbintuition}[why give up expressive power?]
A reasoner answers questions such as ``is class $A$ a subclass of $B$?'' or ``is this
ontology consistent?''. In OWL Full no algorithm can always answer them
(undecidability, as for first-order logic); in OWL DL an algorithm always terminates
with the right answer. Every legal OWL Lite ontology is a legal OWL DL ontology, and
every OWL DL ontology is a legal OWL Full ontology.
\end{nbintuition}

\paragraph{The ontology header.} An OWL document starts by declaring the namespaces it
uses and then describes the ontology itself:
\begin{minted}{xml}
<rdf:RDF
   xmlns:owl ="http://www.w3.org/2002/07/owl#"
   xmlns:rdf ="http://www.w3.org/1999/02/22-rdf-syntax-ns#"
   xmlns:rdfs="http://www.w3.org/2000/01/rdf-schema#"
   xmlns:xsd ="http://www.w3.org/2001/XMLSchema#">
<owl:Ontology rdf:about="">
  <rdfs:comment>An example OWL ontology</rdfs:comment>
  <owl:priorVersion rdf:resource="http://www.mydomain.org/uni-ns-old"/>
  <owl:imports rdf:resource="http://www.mydomain.org/persons"/>
  <rdfs:label>University Ontology</rdfs:label>
</owl:Ontology>
\end{minted}
\texttt{owl:imports} includes other ontologies, whose content is assumed to be part of
the current document. It is \textbf{transitive}: if ontology $A$ imports $B$ and $B$
imports $C$, then $A$ also imports $C$.

\paragraph{Classes.} Classes are defined with \texttt{owl:Class}; further information
can be added to a class later, e.g.\ that it is \textbf{disjoint} from another one (no
individual belongs to both):
\begin{minted}{xml}
<owl:Class rdf:ID="Mammalia">
  <rdfs:subClassOf rdf:resource="#Animal"/>
</owl:Class>
<owl:Class rdf:about="#Mammalia">
  <owl:disjointWith rdf:resource="#Insecta"/>
</owl:Class>
\end{minted}
There are two predefined classes: \texttt{owl:Thing} (the class of everything, the top
of every hierarchy) and \texttt{owl:Nothing} (the empty class).

\paragraph{Properties.} OWL distinguishes two types of property:
\begin{itemize}[nosep]
  \item \textbf{object properties} link individuals to individuals (\emph{married}
        links a person to a person);
  \item \textbf{datatype properties} link individuals to data values (\emph{age} links a
        person to a non-negative integer).
\end{itemize}
\begin{minted}{xml}
<owl:DatatypeProperty rdf:ID="age">
  <rdfs:domain rdf:resource="#Person"/>
  <rdfs:range
  rdf:resource="http://www.w3.org/2001/XMLSchema#nonNegativeInteger"/>
</owl:DatatypeProperty>
<owl:ObjectProperty rdf:ID="married">
  <rdfs:domain rdf:resource="#Person"/>
  <rdfs:range rdf:resource="#Person"/>
</owl:ObjectProperty>
\end{minted}

\begin{nbwarning}[corrections to the slides]
The listings above fix three small slips in the slides: the \texttt{xsd} namespace is
\texttt{XMLSchema} (not \texttt{XLMSchema}); domain and range of OWL properties are
still given with \texttt{rdfs:domain} and \texttt{rdfs:range} (OWL has no
\texttt{owl:domain}/\texttt{owl:range}); and \texttt{rdf:about} on a class defined
elsewhere refers to it with a \texttt{\#}.
\end{nbwarning}

\paragraph{References.} The specifications are on the W3C site: RDF
(\url{http://www.w3.org/RDF/}), RDF Schema (\url{http://www.w3.org/TR/rdf-schema/}) and
OWL (\url{http://www.w3.org/2004/OWL/}).

\subsection{The Semantic Web}
\label{sec:semantic-web}
\slidesrc{Ontologies: 8, 39--48}

\subsubsection{The vision}
\slidesrc{Ontologies: 8, 40--41}

\begin{nbdef}{Semantic Web}
``The \textbf{Semantic Web} is not a separate Web but an extension of the current one,
in which information is given a well-defined meaning, better enabling computers and
people to work in cooperation.'' (Tim Berners-Lee, James Hendler and Ora Lassila,
\emph{Scientific American}, May 2001.)
\end{nbdef}

The same article continues: the Semantic Web will bring structure to the meaningful
content of Web pages, creating an environment where \textbf{software agents} roaming
from page to page can readily carry out sophisticated tasks for users. The goals are
\begin{itemize}[nosep]
  \item automatic \textbf{discovery} of Web resources and Web services;
  \item automatic \textbf{composition} of Web resources and Web services;
  \item automatic \textbf{execution and monitoring} of Web services.
\end{itemize}

\begin{nbintuition}[why the Web needs ontologies]
A page saying ``Dr.\ Rossi, Tuesday 10--12, room 506'' is perfectly clear to a human
and opaque to a program: is it an office hour, a lecture, an appointment? If the page
also states, in RDF, that this is an instance of \emph{OfficeHour} with properties
\emph{teacher}, \emph{day} and \emph{room}, defined in a shared ontology, an agent can
find all office hours of a department, compare them with your calendar and book one.
Ontologies give the shared vocabulary; RDF gives the format of the facts.
\end{nbintuition}

\subsubsection{The Semantic Web stack}
\slidesrc{Ontologies: 42--45}

The Semantic Web technologies are usually drawn as a \textbf{stack of layers}, each
built on the ones below, like the language evolution of
Figure~\ref{fig:languages-evolution}. The stack itself evolved over the years.

\begin{center}
\begin{tikzpicture}[
    x=1.3cm, y=0.55cm]
  % ---- 2003 stack
  \fill[nbgray!25] (0,0) rectangle (1.95,1);
  \node[font=\scriptsize] at ($(0,0)!0.5!(1.95,1)$) {URI};
  \fill[nbgray!25] (2.05,0) rectangle (4,1);
  \node[font=\scriptsize] at ($(2.05,0)!0.5!(4,1)$) {Unicode};
  \fill[nborange!25] (0,1.1) rectangle (1.95,2.1);
  \node[font=\scriptsize] at ($(0,1.1)!0.5!(1.95,2.1)$) {XML};
  \fill[nborange!25] (2.05,1.1) rectangle (4,2.1);
  \node[font=\scriptsize] at ($(2.05,1.1)!0.5!(4,2.1)$) {Namespaces};
  \fill[nborange!45] (0.4,2.2) rectangle (3,3.2);
  \node[font=\scriptsize] at ($(0.4,2.2)!0.5!(3,3.2)$) {RDF M\&S};
  \fill[nborange!45] (0.6,3.3) rectangle (3,4.3);
  \node[font=\scriptsize] at ($(0.6,3.3)!0.5!(3,4.3)$) {RDF Schema};
  \fill[nbblue!20] (0.8,4.4) rectangle (3,5.4);
  \node[font=\scriptsize] at ($(0.8,4.4)!0.5!(3,5.4)$) {Ontology};
  \fill[nbgreen!25] (1.2,5.5) rectangle (3,6.5);
  \node[font=\scriptsize] at ($(1.2,5.5)!0.5!(3,6.5)$) {Rules};
  \fill[nbgreen!25] (1.0,6.6) rectangle (3,7.8);
  \node[font=\scriptsize] at ($(1.0,6.6)!0.5!(3,7.8)$) {Logic framework};
  \fill[nbteal!25] (1.8,7.9) rectangle (3,8.9);
  \node[font=\scriptsize] at ($(1.8,7.9)!0.5!(3,8.9)$) {Proof};
  \fill[nbpurple!20] (2.0,9.0) rectangle (3.9,10);
  \node[font=\scriptsize] at ($(2.0,9.0)!0.5!(3.9,10)$) {Trust};
  \fill[nbred!25] (3.1,2.2) rectangle (3.5,8.9);
  \node[font=\scriptsize, rotate=90] at ($(3.1,2.2)!0.5!(3.5,8.9)$) {Signature};
  \fill[nbred!25] (3.55,2.2) rectangle (3.95,8.9);
  \node[font=\scriptsize, rotate=90] at ($(3.55,2.2)!0.5!(3.95,8.9)$) {Encryption};
  \node[font=\small\bfseries] at (2,-0.8) {2003};
  % ---- 2005 stack
  \begin{scope}[xshift=7cm]
  \fill[nbgray!25] (0,0) rectangle (1.95,1);
  \node[font=\scriptsize] at ($(0,0)!0.5!(1.95,1)$) {URI};
  \fill[nbgray!25] (2.05,0) rectangle (4,1);
  \node[font=\scriptsize] at ($(2.05,0)!0.5!(4,1)$) {Unicode};
  \fill[nborange!25] (0,1.1) rectangle (1.95,2.1);
  \node[font=\scriptsize] at ($(0,1.1)!0.5!(1.95,2.1)$) {XML};
  \fill[nborange!25] (2.05,1.1) rectangle (4,2.1);
  \node[font=\scriptsize] at ($(2.05,1.1)!0.5!(4,2.1)$) {Namespaces};
  \fill[nbred!45] (0,2.2) rectangle (3,3.2);
  \node[font=\scriptsize] at ($(0,2.2)!0.5!(3,3.2)$) {RDF Core};
  \fill[nborange!45] (0.6,3.3) rectangle (3,4.3);
  \node[font=\scriptsize] at ($(0.6,3.3)!0.5!(3,4.3)$) {RDF Schema};
  \fill[nborange!45] (0.6,4.4) rectangle (3,5.4);
  \node[font=\scriptsize] at ($(0.6,4.4)!0.5!(3,5.4)$) {DLP bit of OWL/Rules};
  \fill[nbblue!20] (0.6,5.5) rectangle (1.75,6.5);
  \node[font=\scriptsize] at ($(0.6,5.5)!0.5!(1.75,6.5)$) {OWL};
  \fill[nbgreen!25] (1.85,5.5) rectangle (3,6.5);
  \node[font=\scriptsize] at ($(1.85,5.5)!0.5!(3,6.5)$) {Rules};
  \fill[nbblue!35] (0,3.3) rectangle (0.5,6.5);
  \node[font=\scriptsize, rotate=90] at ($(0,3.3)!0.5!(0.5,6.5)$) {SPARQL};
  \fill[nbgreen!25] (0.8,6.6) rectangle (3,7.8);
  \node[font=\scriptsize] at ($(0.8,6.6)!0.5!(3,7.8)$) {Logic framework};
  \fill[nbteal!25] (1.4,7.9) rectangle (3,8.9);
  \node[font=\scriptsize] at ($(1.4,7.9)!0.5!(3,8.9)$) {Proof};
  \fill[nbpurple!20] (2.0,9.0) rectangle (3.9,10);
  \node[font=\scriptsize] at ($(2.0,9.0)!0.5!(3.9,10)$) {Trust};
  \fill[nbred!25] (3.1,3.3) rectangle (3.5,8.9);
  \node[font=\scriptsize, rotate=90] at ($(3.1,3.3)!0.5!(3.5,8.9)$) {Signature};
  \fill[nbred!25] (3.55,3.3) rectangle (3.95,8.9);
  \node[font=\scriptsize, rotate=90] at ($(3.55,3.3)!0.5!(3.95,8.9)$) {Encryption};
  \node[font=\small\bfseries] at (2,-0.8) {2005};
  \end{scope}
\end{tikzpicture}
\captionof{figure}{The Semantic Web stack according to Tim Berners-Lee, in 2003
(\emph{WWW past \& future}) and in 2005 (\emph{Web for real people}).}
\label{fig:sw-stack}
\end{center}

Reading the 2003 stack from the bottom:
\begin{itemize}[nosep]
  \item \textbf{URI} and \textbf{Unicode}: how to name any resource, how to write any
        character;
  \item \textbf{XML} and \textbf{namespaces}: a common syntax for structured documents,
        and a way to mix vocabularies without name clashes;
  \item \textbf{RDF} (model \& syntax) and \textbf{RDF Schema}: facts as triples, and a
        basic vocabulary of classes and properties
        (Section~\ref{sec:ontology-languages});
  \item \textbf{Ontology} (OWL): rich descriptions of the concepts of a domain;
  \item \textbf{Rules} and \textbf{logic framework}: inference beyond what the ontology
        language can express;
  \item \textbf{Proof} and \textbf{Trust}: explaining why a conclusion holds, and
        deciding whether to believe it (and its sources);
  \item \textbf{Signature} and \textbf{encryption} run alongside all the upper layers:
        who said what, and who can read it.
\end{itemize}

In the 2005 version the ontology layer is split into \textbf{OWL} and \textbf{Rules}
side by side, sharing a common fragment (\emph{DLP}, the part of OWL that can also be
written as rules), and the query language \textbf{SPARQL} appears next to them.

\begin{nbfact}{Stack or two towers?}
Horrocks, Parsia, Patel-Schneider and Hendler (\emph{Semantic Web Architecture: Stack or
Two Towers?}, 2005) argued that a single stack does not really work: rules with
\emph{negation as failure} (NAF, as in Prolog and Datalog) and the open-world logic of
OWL do not combine smoothly. They proposed \textbf{two towers} sharing the bottom layers
(URI, Unicode, XML, namespaces): one with \emph{DLP/Datalog}, rules and NAF, the other
with RDF, RDF Schema, OWL and first-order logic (FOL/SCL), possibly joined at the top by
a common logic framework. As an alternative they also drew a single stack in which the
rule layer (SWRL) sits \emph{above} OWL and below full FOL. The slides illustrate the
debate with Genoa landmarks: the two towers of Porta Soprana for the two-towers view,
the Lanterna (the lighthouse) for the single stack.
\end{nbfact}

\subsubsection{What is the Semantic Web really about?}
\slidesrc{Ontologies: 46}

James Hendler (2009) distinguishes three meanings in which the term is used:
\begin{itemize}[nosep]
  \item \textbf{Linked Data} (or Web of Data): datasets rendered into RDF, using very
        simple ontologies (really more like thesauri or vocabularies than AI languages)
        to provide terms for linking and describing the data;
  \item \textbf{Web 3.0}: sometimes used as a synonym for Semantic Web; it grew up as a
        way of explaining how semantics (especially the simple ontologies above) can be
        integrated with Web 2.0 technologies;
  \item \textbf{Semantic Web}: some people use it to mean all of the above (Hendler
        does).
\end{itemize}

\begin{nbwhy}
Linked Data is where the Semantic Web had its largest practical impact: resources such
as DBpedia and Wikidata are huge RDF knowledge graphs (Section~\ref{sec:what-is-an-ontology}),
and they are used in NLP for entity linking (connecting a name found by NER,
Section~\ref{sec:named-entity-recognition}, to the entity it denotes), question answering
and as background knowledge for language models.
\end{nbwhy}

\subsubsection{Web 4.0 and Web 5.0}
\slidesrc{Ontologies: 47--48}

\begin{itemize}
  \item \textbf{Web 4.0}: its services will be autonomous, proactive, self-learning,
        content-exploring, collaborative and content-generating \textbf{agents}, based on
        fully matured semantic and reasoning technologies as well as AI. They will
        support adaptive content presentation through an intelligent agent; examples are
        services interacting with sensors and implants, natural-language services and
        virtual reality services.
  \item \textbf{Web 5.0}: centred on users' \textbf{control over identity} (credentials
        for all platforms stored on a decentralised platform controlled only by the
        user, with no identity information collected by the Internet) and \textbf{data
        ownership} (complete decentralisation: users choose what to share with
        companies, and store their data with their own blockchain participation).
\end{itemize}
As the slides conclude, people no longer seem very interested in numbering the Web
versions.

\begin{nbsummary}[ontologies and the Semantic Web]
\begin{itemize}[nosep]
  \item \textbf{Ontology}: an explicit specification of a conceptualization (Gruber);
        classes, properties and relations with their meaning and constraints, at the
        ``semantic'' level, closer to first-order logic than database schemas. Main
        motivation: \textbf{semantic interoperability}.
  \item \textbf{Knowledge graph}: interlinked entities and facts described with an
        ontology; new facts can be inferred.
  \item \textbf{Methodology} (Noy \& McGuinness): scope, reuse, terms, classes and
        hierarchy, slots, facets, instances. Tool: \textbf{Protégé}.
  \item \textbf{RDF}: triples subject--predicate--object, forming a graph.
        \textbf{RDFS}: classes, \texttt{subClassOf}, properties with domain and range.
        \textbf{OWL}: richer primitives (disjointness, object vs.\ datatype properties,
        imports); Full (undecidable), DL (efficient reasoning), Lite (simplest).
  \item \textbf{Semantic Web}: the Web with machine-understandable meaning, organised as
        a layered stack (URI/XML $\to$ RDF $\to$ ontologies/rules $\to$ logic, proof,
        trust); in practice, mostly Linked Data.
\end{itemize}
\end{nbsummary}

